\section{Problem Description}
\label{sec:problem}

A bounded planar workspace is considered with immovable structures, a set of
movable obstacles, and a small team of robots operating cooperatively. Rather
than simply navigating through the scene, the robots actively reconfigure it by
pushing selected obstacles. The task is to enable an external vehicle to travel
from a start location to a goal location through a sufficiently wide corridor.
Formally, a solution is valid only if there exists a path for the vehicle whose
clearance is no smaller than $W$. This requirement defines the task at the
level of corridor existence rather than robot-only reachability.

A clearing action is not treated as an abstract obstacle displacement. Instead,
each action is represented as a pushing mode, which specifies robot-to-contact
assignments, contact locations on obstacle boundaries, and the applied forces
or body wrench used to induce motion. The resulting planning problem is hybrid
in nature. The system must decide which blocked gap or obstacle to clear at the
discrete level, while simultaneously determining how to realize that action
through contact-feasible pushing at the continuous level. Collaborative path
clearing is therefore a contact-rich decision problem rather than a purely
geometric graph-search problem.
